\section[The $\cR$-function as a special case of the $R$-function]{The $\boldsymbol{\cR}$-function as a special case of the $\boldsymbol{R}$-function}\label{section2}

\subsection[The functions $\cR$ and $\cR_r$]{The functions $\boldsymbol{\cR}$ and $\boldsymbol{\cR_r}$}\label{section2.1}

The $R$-function \eqref{Rdef} is def\/ined as a contour integral over a variable $z$, with the $z$-dependence of the integrand encoded in a product of eight hyperbolic gamma function factors.
(See Appendix~\ref{appendixA} for a review of the relevant features of the hyperbolic gamma function.)
Specif\/ically,
with suitable restrictions on the eight variables, the $R$-function is given by
\begin{gather}\label{R}
R({\bf c};v,\hat{v})=\sqrt{\frac{\alpha}{2\pi}}\int_{\cal
C}F(c_0;v,z)K({\bf c};z)F(\hat{c}_0;\hat{v},z)dz.
\end{gather}
Here we have
\begin{gather}
\hat{c}_0\equiv (c_0+c_1+c_2+c_3)/2,
\\
\label{F}
F(d;y,z)\equiv
\frac{G(z\pm y+id-ia)}{G(\pm y+id-ia)},
\end{gather}
(with $f(w\pm y)$ denoting $f(w+y)f(w-y)$), and $K$ is given by
\begin{gather}\label{K}
K({\bf c};z)\equiv
\frac{1}{G(z+ia)}\prod_{j=1}^3\frac{G(is_j)}{G(z+is_j)},
\end{gather}
with new parameters
\begin{gather}\label{defsj}
s_1\equiv c_0+c_1-a_{-}/2,\qquad s_2\equiv c_0+c_2-a_{+}/2,\qquad s_3\equiv
c_0+c_3.
\end{gather}
Also, recall $a$ and $\alpha$ are def\/ined by~\eqref{aconv}.

We do not need the def\/inition of the contour $\cC$ for general variable choices (this is discussed in I and Section~4 of the survey~\cite{Rsurv}); instead we presently def\/ine $\cC$ for the cases at issue. For the special ${\bf c}$-choice in \eqref{cRdef} we can use the duplication formula~\eqref{dupl} to obtain
\begin{gather}
K((b,0,0,0);z)=\cK(b;z),
\end{gather}
where
\begin{gather}
\cK(b;z)\equiv \frac{1}{G(z+ia)}\frac{G(2ib-ia)}{G(ib-ia)}\frac{G(z+ib-ia)}{G(2z+2ib-ia)}.
\end{gather}
Using also the ref\/lection equation~\eqref{refl}
we deduce that $\cR$ is given by
\begin{gather}
\cR(b;x,y)     =     \left(\frac{\alpha}{2\pi}\right)^{1/2}\frac{G(2ib-ia)}{G(ib-ia)}G(\pm x-ib+ia)G(\pm y/2-ib/2+ia)
\nonumber \\
\phantom{\cR(b;x,y)     =}{}\times \int_{\cal
C}\frac{G(z\pm x+ib-ia)G(z+ib-ia)G(z\pm y/2+ib/2-ia)}{G(z+ia)G(2z+2ib-ia)}dz.\label{cRcont}
\end{gather}

For the variable choice that is most relevant for Hilbert space purposes, namely,
\begin{gather}\label{Hpar}
\aaa,b,x,y>0,
\end{gather}
the contour $\cC$ may be chosen equal to the real line in the $z$-plane, indented downwards near 0 so as to avoid a pole of $\cK(b;z)$. From~\eqref{Gpo}--\eqref{Gze} it follows that the poles of $\cK(b;z)$ are located on the imaginary axis at
\begin{gather}\label{Kpoles}
z-z_{kl}=0,\qquad z-z_{kl}+ib=ia_{+}/2, \, ia_{-}/2, \, 0,\qquad k,l\in\N.
\end{gather}
Thus they are above the contour, whereas the remaining $z$-poles of the integrand at
\begin{gather}
z+z_{kl}=\pm x-ib,\  \pm y/2-ib/2,\qquad  k,l\in\N,
\end{gather}
are below $\cC$.

From the above representation it is immediate that $\cR$ is symmetric under the interchange of the parameters $a_{+}$ and $a_{-}$:
\begin{gather}\label{Rmodinv}
\cR(a_{+},a_{-},b;x,y)=\cR(a_{-},a_{+},b;x,y).
\end{gather}
It is not at all clear, however, that $\cR$ is also symmetric under the interchange of the positions~$x$ and~$y$:
\begin{gather}\label{sd}
\cR(a_{+},a_{-},b;x,y)=\cR(a_{+},a_{-},b;y,x).
\end{gather}
 This self-duality feature follows in particular from a second relation between $\cR$ and $R$, namely,
\begin{gather}\label{cR2}
\cR(b;x,y)=R((b,b,b,b)/2;x/2,y).
\end{gather}
(This is equation~(4.8) in~\cite{quad}.) Indeed, this second ${\bf c}$-choice yields the same function $\cK(b;z)$ as the f\/irst one, so that substitution of~\eqref{R} (with the same contour $\cC$) now yields \eqref{cRcont} with $x$ and $y$ interchanged on the r.h.s.

There are two more ${\bf c}$-choices that lead from $R$ to $\cR$, namely, $(b,0,b,0)$ and $(b,b,0,0)$. Specif\/ically, from equations~(4.6) and (4.7) in~\cite{quad} we have
\begin{gather}\label{cR3}
\cR(\aaa,b;x,y)=R(a_{+},2a_{-},(b,0,b,0);x,y),
\\
\label{cR4}
\cR(\aaa,b;x,y)=R(2a_{-},a_{+},(b,b,0,0);x,y).
\end{gather}
From the def\/inition of the $R$-function we then obtain alternative integral representations for the $\cR$-function from which the self-duality property \eqref{sd} is manifest. (Indeed, since we have $c_0=\hat{c}_0=b$ for these two choices, the integrand is invariant under the interchange of $x$ and $y$.) On the other hand, the modular invariance property \eqref{Rmodinv} is not at all clear, since the integral representations involve the hyperbolic gamma function
 with $a_{-}$ replaced by $2a_{-}$. Using~\eqref{G2},
 they can be re-expressed in terms of the modular invariant function $G(\aaa;z)$. However, the resulting integrand is then still not modular invariant. Since it seems not to simplify and does not look
illuminating, we do not detail it any further.

The analyticity properties of the $R$-function are known in great detail from~Theorem~2.2 in I, cf.~also Section~4 in the survey~\cite{Rsurv}. Combining this theorem with the def\/inition~\eqref{cRdef} of $\cR$ and its self-duality property~\eqref{sd}, we deduce in particular that $\cR$ extends from the intervals \eqref{Hpar} to a meromorphic function in~$b$,~$x$ and~$y$, whose poles in~$x$ and~$y$ can only occur at the locations
\begin{gather}\label{cRpoles}
\pm z=2ia-ib+z_{kl},\qquad  z=x,y,\qquad k,l\in \N.
\end{gather}

We proceed to list further consequences of the $R$-function theory for $\cR$. Two features that are clear from each of the above integral representations are evenness and scale invariance (given scale invariance of $G$):
\begin{gather}\label{even}
\cR (b;x,y)=\cR(b;\de x, \de'y),\qquad \de,\de'=+,-,
\\
\label{scale}
\cR(\aaa,b;x,y)=\cR(\lambda  a_{+},\lambda a_{-}, \lambda b;\lambda x,\lambda y),\qquad \lambda>0.
\end{gather}
A less obvious feature is the explicit evaluation
\begin{gather}\label{cRnorm}
\cR(b;x,ib)=1.
\end{gather}
It follows from the formula
\begin{gather}
R({\bf c};v,i\hat{c}_0)=1,
\end{gather}
(cf.~equation~(3.26) in I or Section~6 in~\cite{Rsurv}), by using any of the four relations \eqref{cRdef}, \eqref{cR2}, \eqref{cR3}, \eqref{cR4}.
Def\/ining next
\begin{gather}
\cR_n(x)\equiv \cR(b;x,y_n),\qquad y_n\equiv ib+ina_{-},\qquad n\in\N,
\end{gather}
the eigenvalue A$\De$E (analytic dif\/ference equation) for $A_{+}(b;y)$ entails
\begin{gather}
 \frac{s_{+}(y_n-ib)}{s_{+}(y_n)}\cR_{n-1}(x)+\frac{s_{+}(y_n+ib)}{s_{+}(y_n)}\cR_{n+1}(x)=2c_{+}(x)\cR_n(x).
\end{gather}
In view of \eqref{cRnorm}, it follows from this that $\cR_n$ is of the form
\begin{gather}\label{qGeg}
\cR_n(x)=P_n(c_{+}(x)),
\end{gather}
where $P_n(z)$ is a polynomial in $z$ of degree $n$ and parity $(-)^n$. The relation of these polynomials to the $q$-Gegenbauer polynomials and to the  Askey--Wilson polynomials associated with the four relevant ${\bf c}$-choices is detailed at the end of Section~4 of~\cite{quad}.

The renormalized $\cR$-function~$\cR_r$ given by~\eqref{cRr}
is the counterpart of the renormalized $R$-function $R_r$ obtained from \eqref{R} by omitting the product $\prod_jG(is_j)$ in $K$, cf.~\eqref{K}. (To see this, use \eqref{refl} and \eqref{dupl}.) Clearly, it shares the features \eqref{Rmodinv}, \eqref{sd}, \eqref{even} and \eqref{scale} of the $\cR$-function, whereas \eqref{cRnorm} is replaced by
\begin{gather}\label{cRrnorm}
\cR_r(b;x,ib)=\frac{G(ib-ia)}{G(2ib-ia)}.
\end{gather}
The renormalizing factor in the function~$\cR_r$ ensures that it has no poles that are independent of $x$ and $y$, cf.~Theorem~2.2 in~I. More precisely, $\cR_r(\aaa,b;x,y)$ extends to a function that is meromorphic in the domain{\samepage
\begin{gather}\label{Dp}
D_+\equiv \big\{ (\aaa,b,x,y)\in \C^5 \mid \re a_+>0,\ \re a_->0\big\},
\end{gather}
and whose poles can only occur at the locations~\eqref{cRpoles}.}

It is not obvious, but true that for the special $b$-choices $b_{mn}$~\eqref{bmn}
we have an equality
\begin{gather}\label{cRM}
\cR_r(b_{mn};x,y)=M(b_{mn};x,y)+M(b_{mn};-x,y),
\end{gather}
where $M(b_{mn};x,y)$ is the function def\/ined at the end of Section~III in our paper~\cite{GLFII}. Therefore, $\cR_r(b;x,y)$ is the continuous (indeed, real-analytic) interpolation to arbitrary $b\in\R$ of  the function given by equation~(3.74) in~\cite{GLFII}, which is def\/ined only for the $b$-values~$b_{mn}$. (Note the latter are dense in $\R$ when the ratio $a_{+}/a_{-}$ is irrational.)

It will not cause surprise that in the free case we have
\begin{gather}\label{M00}
M(b_{00};x,y)=\exp(i\alpha xy/2).
\end{gather}
It would take us too far af\/ield, however, to detail all of the functions $M(b_{mn};x,y)$, $m,n\in\Z$. For our purposes it is enough to specify their general structure: They are `elementary' in the sense that they can be written
\begin{gather}\label{Mform}
M(b_{mn};x,y)=\exp(i\alpha xy/2)R_{mn}(e_{+}(x),e_{-}(x),e_{+}(y),e_{-}(y)),
\end{gather}
where $R_{mn}$ is a rational function of its four arguments, cf.~Section~III in~\cite{GLFII}. In Subsection~\ref{section2.4} we deduce this structure in another way (namely, by exploiting parameter shifts). Moreover, for the case where $m$ and $n$ are not both positive or both non-positive, this structure can be understood from the novel Fourier transform representations~\eqref{Rii}--\eqref{Rv}, cf.~Subsection~\ref{section4.1}.

We postpone the proof of the equality assertion \eqref{cRM} to Subsection~\ref{section2.3}. An ingredient of this proof is the asymptotic behavior of $\cR_r(b;x,y)$ as $x$ goes to $\infty$, and this is most easily obtained as a corollary of the asymptotics of a closely related function~$\rE(b;x,y)$, def\/ined by~\eqref{EcE}.

\subsection[The functions $\rE$ and $\rF$]{The functions $\boldsymbol{\rE}$ and $\boldsymbol{\rF}$}\label{section2.2}

The function $\rE(b;x,y)$ can be viewed as a specialization of the function denoted $\cE(\gamma;v,\hat{v})$ in II and~\cite{Rsurv}. The relation between $\gamma$ and ${\bf c}$ reads
\begin{gather}\label{gamc}
\gamma(\aaa,{\bf c})=(c_0-a,c_1-a_{-}/2,c_2-a_{+}/2,c_3).
\end{gather}
In particular, the `free' case ${\bf c}={\bf 0}$ yields
\begin{gather}\label{gamf}
\gamma_{\rm f}\equiv (-a,-a_{-}/2,-a_{+}/2,0).
\end{gather}
The switch from ${\bf c}$ to $\gamma$ is crucial for uncovering further symmetries: The function $\cE(\gamma;v,\hat{v})$ is invariant under $D_4$ transformations on $\gamma$ (i.e., permutations and even sign changes), cf.~II. (In~\cite{bust} this $D_4$ symmetry has been reobtained in a quite dif\/ferent way.) It is def\/ined by
\begin{gather}
\cE(\gamma;v,\hat{v})=\frac{\chi(\gamma)}{c(\gamma;v)c(\hat{\gamma};\hat{v})}R_r(\gamma;v,\hat{v}).
\end{gather}
Here, the generalized ($BC_1$) Harish-Chandra $c$-function is given by
\begin{gather}\label{c}
c(\gamma;v)\equiv \frac{1}{G(2v+ia)}\prod_{\mu =0}^3G(v-i\gamma_{\mu}),
\end{gather}
the dual of $\gamma$ by
\begin{gather}\label{dualgam}
\hat{\gamma}\equiv J\gamma,\qquad
J\equiv \frac{1}{2} \left( \begin{array}{rrrr}
1  &  1  &  1  &  1  \\
1  &  1  &  -1  &  -1  \\
1  &  -1  &  1  &  -1  \\
1  &  -1  &  -1  &  1
\end{array} \right) ,
\end{gather}
and the constant by
\begin{gather}\label{chigam}
\chi(\gamma) \equiv \exp\big(i\alpha [\gamma\cdot \gamma /4-\big(a_{+}^2+a_{-}^2+a_{+}a_{-}\big)/8]\big),\qquad
\alpha =2\pi/a_{+}a_{-}.
\end{gather}

Denoting the $\gamma$-vectors corresponding to the two one-parameter families
\begin{gather}
{\bf c}=(b,0,0,0), (b,b,b,b)/2,
\end{gather}
by $\gamma^{(1)}$,  $\gamma^{(2)}$, it is easy to check that
\begin{gather}\label{J12}
J\gamma^{(1)}=\gamma^{(2)}.
\end{gather}
Def\/ining
\begin{gather}\label{EcE}
\rE(b;x,y)\equiv \cE\big(\gamma^{(1)};x,y/2\big),
\end{gather}
a straightforward calculation (using the duplication formula \eqref{dupl}) yields
\begin{gather}\label{ER}
\rE(b;x,y)=\frac{\phi(b)}{c(b;x)c(b;y)}\cR_r(b;x,y),
\end{gather}
where we have introduced the constant
\begin{gather}\label{phib}
\phi(b)\equiv \exp(i\alpha b(b-2a)/4),
\end{gather}
and generalized ($A_1$) Harish-Chandra $c$-function
\begin{gather}\label{cb}
c(b;z)\equiv \frac{G(z+ia-ib)}{G(z+ia)}.
\end{gather}
Recalling \eqref{cR2} and using \eqref{J12}, it readily follows that we also have
\begin{gather}
\rE(b;x,y)=\cE\big(\gamma^{(2)};x/2,y\big).
\end{gather}

It involves more work to obtain the relations between $\rE$ and $\cE$ corresponding to \eqref{cR3} and~\eqref{cR4}. Setting
\begin{gather}
\gamma^{(3)}\equiv \gamma(a_{+},2a_{-},(b,0,b,0)),\qquad \gamma^{(4)}\equiv \gamma(2a_{-},a_{+},(b,b,0,0)),
\end{gather}
these are given by
\begin{gather}
\rE(\aaa,b;x,y)=\cE\big(a_{+},2a_{-},\gamma^{(3)};x,y\big)=\cE\big(2a_{-},a_{+},\gamma^{(4)};x,y\big).
\end{gather}
(These formulas amount to special cases of the doubling identity for the $\cE$-function obtained in Section~6 of~\cite{quad}.)

The relation \eqref{ER} between $\rE$ and $\cR_r$ yields a similarity transformation turning the A$\De$Os \eqref{4A} into $\cA_{\pm}(b;x)$, $\cA_{\pm}(b;y)$,
where
\begin{gather}\label{cAdef}
\cA_{\de}(b;z)\equiv c(b;z)^{-1}A_{\de}(b;z)c(b;z)=T^z_{ia_{-\de}}+V_{\de}(b;z)T^z_{-ia_{-\de}},\qquad \de=+,-,
\\
V_{\de}(b;z)\equiv \frac{s_{\de}(z+ib)s_{\de}(z-ib +ia_{-\de})}{s_{\de}(z)s_{\de}(z+ia_{-\de})}.
\end{gather}
From this it is easy to verify that these A$\De$Os are formally self-adjoint operators on $L^2(\R)$ (by contrast to the A$\De$Os \eqref{4A}), and that they are invariant under the transformation
\begin{gather}\label{bsym}
b \mapsto a_{+}+a_{-}-b.
\end{gather}
It is not obvious, but true that we also have
\begin{gather}\label{Esym}
\rE(b;x,y)=\rE(a_{+}+a_{-}-b;x,y).
\end{gather}
This symmetry property can be derived from \eqref{EcE} and the $D_4$ invariance of the $\cE$-function: We have
\begin{gather}\label{gam1}
\gamma^{(1)}=(b-a,-a_{-}/2,-a_{+}/2,0),
\end{gather}
so the map $b\mapsto 2a-b$ yields a $\gamma$-vector related to $\gamma^{(1)}$ by a sign f\/lip of the f\/irst and last component.

Combining~\eqref{ER} and~\eqref{cb} with the analyticity features of $\cR$ (cf.~the paragraph con\-taining~\eqref{cRpoles}), we deduce that $\rE(b;x,y)$ is meromorphic in $b$, $x$ and $y$, with $b$-independent pole locations
\begin{gather}
z=-2ia-z_{kl},\qquad z=x,y,\qquad k,l\in\N,
\end{gather}
corresponding to the factor $G(x+ia)G(y+ia)$, and $b$-dependent poles at
\begin{gather}
z=ib+z_{kl},\qquad z=-ib+2ia+z_{kl},\qquad z=x,y,\qquad k,l\in\N.
\end{gather}

The main disadvantage of the function $\rE(b;x,y)$ compared to the $\cR_r$-function is that it is not even in $x$ and $y$, since the $c$-functions in \eqref{ER} are not even. Instead, it satisf\/ies
\begin{gather}\label{Erefl}
\rE(b;-x,y)=-u(b;x)\rE(b;x,y),
\end{gather}
where
\begin{gather}\label{u}
u(b;z)\equiv -c(b;z)/c(b;-z)=-G(z\pm (ia-ib))/G(z\pm ia).
\end{gather}
On the other hand, $\rE$ inherits all other important properties of $\cR_r$, and is the simplest function to use for Hilbert space purposes. In particular, it has the `unitary asymptotics'
\begin{gather}\label{Eas}
\rE(b;x,y)\sim \exp(i\alpha xy/2)- u(b;-y)\exp(-i\alpha xy/2),\\
 b\in\R,\qquad y\in(0,\infty), \qquad x\to \infty,\nonumber
\end{gather}
cf.~Theorem~1.2 in~II. Here, the $u$-function
encodes the scattering associated with the A$\De$Os $\cA_{\pm}(b;x)$, reinterpreted as commuting self-adjoint operators on the Hilbert space $L^2((0,\infty),dx)$. More precisely, using corresponding results on the $\cE$-function from~II and~III, it follows that the generalized Fourier
transform
\begin{gather}\label{cF}
\cF
  : \ \cC\equiv C_0^{\infty}((0,\infty))\subset \hat{\cH}\equiv L^2((0,\infty),dy) \to\cH
\equiv L^2((0,\infty),dx),
\end{gather}
def\/ined by
\begin{gather}\label{cFE}
(\cF\psi)(x)\equiv\left(\frac{\alpha}{4\pi}\right)^{1/2}\int_0^{\infty}\rE(b;x,y)\psi(y)dy,\qquad
\psi\in\cC,
\end{gather}
extends to a unitary operator, provided the coupling $b$ is suitably restricted. Specif\/ically, it suf\/f\/ices to require
\begin{gather}\label{bres}
b\in [0,a_{+}+a_{-}].
\end{gather}
The self-adjointness of the operators $\hat{\cA}_{\pm}(b)$ on $\cH$ associated to the A$\De$Os $\cA_{\pm}(b;x)$ for $b$ in this interval can then be easily understood from the unitarity of $\cF$: they are the pullbacks to $\cH$ under $\cF$ of the self-adjoint operators of multiplication by $2c_{\pm}(y)$ on $\hat{\cH}$.

As already mentioned, these statements follow from II and III by specialization, but it may help to look f\/irst at Section~9 in the survey~\cite{Rsurv}. Starting from the representation~\eqref{EcE}, the
vector $\gamma^{(1)}$ belongs to the polytope $P$ given by equation~(9.2) in~\cite{Rsurv}, provided $b\in (0,2a)$. Therefore, the transform associated with $\cE(\gamma^{(1)};v,\hat{v})$ (def\/ined by equation~(9.4)) is an isometry. As a consequence, the transform~\eqref{cFE} is an isometry. (The normalization factor in~\eqref{cFE} dif\/fers from that in~equation~(9.4) in~\cite{Rsurv} to accommodate the scale factor 1/2 in the $y$-dependence of~$\cE$ in~\eqref{EcE}.) Isometry of the inverse transform is then clear from the self-duality of $\rE(b;x,y)$. Next, for $b=0$ we have the identity
\begin{gather}\label{ER0}
\rE(0;x,y)=\cR_r(0;x,y)=2\cos(\alpha xy/2),
\end{gather}
cf.~\eqref{ER}--\eqref{cb} and~equation~(7.33) in~\cite{Rsurv} (also, note that for $b=0$ we have $\gamma^{(1)}=\gamma_{\rm f}$, cf.~\eqref{gam1} and~\eqref{gamf}). In view of the symmetry~\eqref{Esym},
it follows that $\cF$ amounts to the cosine transform for $b=0$ and $b=2a$, so these transforms are unitary as well. More generally, we obtain a family of unitary operators
\begin{gather}
\cF(\aaa,b),\qquad (\aaa,b)\in \Pi_u\equiv(0,\infty)^2\times[0,a_{+}+a_{-}],
\end{gather}
which is strongly continuous on the parameter set $\Pi_u$ and satisf\/ies
\begin{gather}\label{cFsym}
\cF(\aaa, b)=\cF(\aaa, a_{+}+a_{-}-b),
\end{gather}
cf.~Theorem~3.3 in III.

The $G$-function asymptotics~\eqref{Gas}
entails that the $c$-function~\eqref{cb} has asymptotics
\begin{gather}\label{cas}
c(b;x)\sim \phi(b)^{\pm 1}\exp(\mp \alpha bx/2),\qquad \re (x)\to\pm\infty,
\end{gather}
with $\phi(b)$ given by~\eqref{phib}. Hence the
$u$-function~\eqref{u} has asymptotics
\begin{gather}\label{uas}
u(b;x)\sim -\phi(b)^{\pm 2},\qquad \re ( x)\to  \pm\infty.
\end{gather}
Also, the ref\/lection equation~\eqref{refl} and the complex conjugation relation~\eqref{Gcon}
entail
\begin{gather}
u(b;-x)u(b;x)=1,\qquad |u(b;x)|=1,\qquad b,x\in\R.
\end{gather}
Thus, if we set
\begin{gather}\label{Fdef}
\rF(b;x,y)\equiv \phi(b)^{-1}(-u(b;x))^{1/2}(-u(b;y))^{1/2}\rE(b;x,y),\qquad b,x,y>0,
\end{gather}
(with the square root phase factors reducing to 1 for $b=0$), then $\rF$ has asymptotics
\begin{gather}\label{Fas}
\rF(b;x,y)\sim [-u(b;y)]^{1/2}\exp(i\alpha xy/2) +[-u(b;y]^{-1/2}\exp(-i\alpha xy/2),\qquad x\to \infty,
\end{gather}
and if we replace $\rE$ by $\rF$ in the above unitary transform~\eqref{cFE}, we retain unitarity.

Introducing the weight function
\begin{gather}\label{w}
w(b;x)\equiv 1/c(b;\pm x)=G(\pm x +ia)/G(\pm x +ia-ib),
\end{gather}
we have
\begin{gather}
w(b;x)>0, \qquad b,x>0,
\end{gather}
and we can also write $\rF$ in terms of $\cR_r$ as
\begin{gather}\label{FR}
\rF(b;x,y) =w(b;x)^{1/2}w(b;y)^{1/2}\cR_r(b;x,y),\qquad b,x,y>0,
\end{gather}
with the positive square roots understood.
Hence, $\rF$ is a joint eigenfunction of the A$\De$Os $H_{\pm}(b;x)$ and
$H_{\pm}(b;y)$
with eigenvalues~\eqref{4ev}, where
\begin{gather} 
H_{\de}(b;z)     \equiv    w(b;z)^{1/2} A_{\de}(b;z)w(b;z)^{-1/2}
\nonumber \\
\label{H} \phantom{H_{\de}(b;z)}{}    =
\sum_{\tau=+,-}\left( \frac{s_{\de}(z-\tau ib)}{s_{\de}(z)}\right)^{1/2}T^z_{\tau ia_{-\de}}\left(\frac{s_{\de}(z+\tau ib)}{s_{\de}(z)}\right)^{1/2}.
\end{gather}

From the $G$-A$\De$Es~\eqref{Gades}
we deduce
\begin{gather}\label{wwr}
w(b;x) =4s_{+}(x)s_{-}(x)w_r(b;x), \qquad w_r(b;x)\equiv G(\pm x -ia+ib).
\end{gather}
For $b\in (0,2a)$ the reduced weight function $w_r(b;x)$ is positive for all real $x$, and since it is also even, its positive square root for $x>0$ has a real-analytic extension to an even positive function on all of~$\R$. By contrast, it is clear from~\eqref{wwr} that $w(b;x)^{1/2}$, $x>0$, extends to an odd real-analytic function on~$\R$. As a consequence, one can also view the transform associated with $\rF(b;x,y)$, $b\in (0,2a)$, as a unitary transform from the odd subspace of $L^2(\R,dy)$ onto the odd subspace of $L^2(\R,dx)$. This is the viewpoint taken in~\cite{Hilb}, where we studied this transform (among other ones) for the special $b$-values $Na_{+}$ with $N\in\N^{*}$. As shown there, for $b>2a$ unitarity and self-adjointness generically break down in a way that can be understood in great detail.

To be sure, the precise connection between the above functions~$\rF$ and~$\cR_r$ and the  functions~$F_r$ and~$E_r$ from~\cite{Hilb} is not clear at face value. But the latter are derived from the functions $M((N+1)a_{+};x,y)$ of~\cite{GLFII}, as specif\/ied below equation~(1.42) in~\cite{Hilb}, so this connection is encoded in the identities~\eqref{cRM} for the special cases $(m,n)=(N+1,0)$, $N\in\N$.

\subsection{The identities (\ref{cRM}) and their consequences}\label{section2.3}

We proceed to prove the general identities~\eqref{cRM}.
Our reasoning involves in particular a comparison of the behavior for $x\to\infty$ of the functions on the l.h.s.\ and r.h.s. For $\cR_r(b;x,y)$ this asymptotics easily follows upon combining~\eqref{ER}, \eqref{Eas} and \eqref{cas}:
\begin{gather}\label{Rras}
\cR_r(b;x,y)\sim \exp(-\alpha bx/2)\sum_{\tau=+,-}c(b;\tau y)\exp(\tau i\alpha xy/2),\\
 b\in\R,\qquad y>0,\qquad x\to\infty.\nonumber
\end{gather}

{\sloppy Next, we consider the functions $M(b_{mn};\pm x,y)$. To begin with, they are eigenfunctions of the four A$\De$Os~\eqref{4A} (where $b=b_{mn}$) with eigenvalues~\eqref{4ev}, cf.~Theorem~II.3 in~\cite{GLFII}. Their `ele\-men\-ta\-ry' form~\eqref{Mform} follows from equations~(3.65)--(3.68) in~\cite{GLFII}. The function $K_{N_{+},N_{-}}(\aaa;x,y)$ occurring in these formulas is specif\/ied in equation~(3.2), with~$S_{N_{\de}}$ given by equation~(2.21). In turn, the coef\/f\/icients in equation~(2.21) are def\/ined via equations~(2.2)--(2.5) in~\cite{GLFII}. (See also Subsection~\ref{section4.1} for more information on these special cases.) It is straightforward to obtain the asymptotics for $\re(x)\to \infty$ from these explicit formulas. Specif\/ically, this yields
\begin{gather}\label{Mas}
M(b_{mn};\pm x,y)=\exp (- \alpha b_{mn}x/2)c(b_{mn};\pm y)e^{ \pm i\alpha xy/2}\big[1+O\big(e^{-\rho \re(x)}\big)\big],\\
 y>0,\qquad \re(x)\to \infty.\nonumber
\end{gather}
The decay rate $\rho$ is the minimum of the two numbers $2\pi /a_{\pm}$, and the implied constant can be chosen uniform for $\im (x)$ varying over $\R$.

}

Comparing \eqref{Rras} and \eqref{Mas}, it follows that the functions on the l.h.s.\ and r.h.s.\ of~\eqref{cRM} have the same asymptotics for $x\to \infty$. It therefore suf\/f\/ices to prove that for f\/ixed $\aaa,y>0$ and $m,n\in\Z$, they must be proportional as functions of $x$. Moreover, we may as well assume $a_{+}/a_{-}$ is irrational, since equality for this case entails equality for all $\aaa>0$. (Indeed, the functions $M(b_{mn};\pm x,y)$ are manifestly real-analytic in $a_{+}$ and $a_{-}$ for $\aaa >0$, and this real-analyticity property is also valid for $\cR_r(b;x,y)$, cf.~I.)

The key consequence of the irrationality assumption is that the vector space of meromorphic joint solutions $f(x)$ to the A$\De$Es
\begin{gather}\label{joint}
A_{\pm}( ma_{+}+na_{-};x)f(x)=2c_{\pm}(y)f(x),\qquad  m,n\in\Z,\qquad y>0,\qquad  a_{+}/a_{-}\notin \Q,
\end{gather}
is two-dimensional.
To explain why this is so, we f\/irst note that the functions $M(b_{mn};\pm x,y)$ are independent solutions to~\eqref{joint}, their independence already being clear from their general form~\eqref{Mform}. Moreover, it follows from their uniform asymptotics~\eqref{Mas} that there exists a~po\-si\-ti\-ve number $\Lambda$, depending on the f\/ixed variables $a_+$, $a_-$, $m$, $n$ and $y$, but not on $x$, such that in the half plane $\re(x)>\Lambda$ both functions are zero-free, and satisfy
\begin{gather}\label{qas}
\lim_{\im(x)\to\infty}M(b_{mn};x,y)/M(b_{mn};-x,y)=0,\qquad \re(x)\in (c_{-},c_{+})\subset [\Lambda, \infty).
\end{gather}

We are now in the position to invoke a result from Section~1 in~\cite{side}, to the ef\/fect that the above  suf\/f\/ices for any joint meromorphic solution $f(x)$ of~\eqref{joint} to be a linear combination of the two functions $M(b_{mn};\pm x,y)$.  Since $\cR_r(b_{mn};x,y)$ is an even meromorphic joint solution, the functions on the l.h.s.\ and r.h.s.\ of~\eqref{cRM} are proportional, so their equality now follows. In particular, for the free case $m=n=0$ we recover the identity \eqref{ER0} from~\eqref{cRM}--\eqref{M00}. Moreover, taking $y=ib_{mn}$ in~\eqref{cRM}, we can invoke~\eqref{cRrnorm} to deduce the corollary
\begin{gather}\label{Mid}
M(b_{mn};x,ib_{mn})+M(b_{mn};-x,ib_{mn})=\frac{G(i(m-1/2)a_{+}+i(n-1/2)a_{-})}{G(i(2m-1/2)a_{+}+i(2n-1/2)a_{-})}.
\end{gather}
Using the $G$-A$\De$Es~\eqref{Gades}, the r.h.s.\ can be rewritten in terms of sine-functions. For the special case $m=N+1$, $n=0$, the resulting identity amounts to equation~(2.78) in~\cite{GLFII}, cf.~also Subsection~\ref{section4.1}.

We would like to add in passing that it is very plausible that~\eqref{qas} is not necessary for two-dimensionality. Indeed, denoting by $\cP_c$ the f\/ield of meromorphic functions with period $c\in\C^{*}$, any third independent joint meromorphic solution would have to be both of the form
\begin{gather}\label{form1}
f(x)=p_1(x)M(b_{mn};x,y)+p_2(x)M(b_{mn};-x,y),\qquad p_1,p_2\in\cP_{ia_{+}},
\end{gather}
and of the form
\begin{gather}\label{form2}
f(x)=q_1(x)M(b_{mn};x,y)+q_2(x)M(b_{mn};-x,y),\qquad  q_1,q_2\in\cP_{ia_{-}}.
\end{gather}
Since the intersection of the f\/ields $\cP_{ia_{+}}$ and $\cP_{ia_{-}}$ reduces
 to the constants when $a_{+}/a_{-}$ is irrational, we expect (but are unable to prove) that this simultaneous representation should lead to a contradiction without appealing to~\eqref{qas}.

Now that we have proved~\eqref{cRM},  it follows that the function $\cR_r(\aaa,b;x,y)$, which is real-analytic on the parameter set
\begin{gather}\label{Pi}
\Pi\equiv \big\{(\aaa,b)\in (0,\infty)^2\times\R\big\},
\end{gather}
 is the continuous interpolation of the functions on the r.h.s.\ of~\eqref{cRM}, which are only def\/ined for the dense subset of `elementary' parameters
\begin{gather}\label{elem}
\Pi_{\rm el}\equiv \{ (\aaa,b)\in\Pi \mid b=ma_{+}+na_{-},\ m,n\in\Z\}.
\end{gather}
The natural question whether another linear combination of $M(b_{mn};\pm x,y)$ that is independent from the even one admits a continuous interpolation as well remains open. In this connection we should point out  that our reasoning at the end of~Section~3 of~\cite{koor} renders this extremely unlikely, but is not conclusive. Indeed, we cannot rule out that the sequence of functions $Q_{-}$ given by equation~(3.15) in~\cite{koor}, with $N_{+}\in\N$, gives rise to an inf\/inity of distinct limits $L_{-}$, corresponding to distinct subsequences. (This oversight is of no consequence for the later sections in~\cite{koor}.) For the same reason, the analogous assertion about the $R$-function, made at the end of~\cite{Rsurv}, has not been completely proved.

Before turning to parameter shifts, we derive  a non-obvious reality feature of $\cR_r$ from the  relations~\eqref{cRM}, namely,
\begin{gather}\label{Rreal}
\overline{\cR_r}(\aaa,b;x,y)=\cR_r(\aaa,b;x,y),\qquad \forall\, (\aaa,b,x,y)\in\Pi\times\R^2.
\end{gather}
The point is that from the explicit formulas for the functions $M$ it is apparent that we have
\begin{gather}
\overline{M}(b_{mn};x,y)=M(b_{mn};-x,y), \qquad x,y\in\R,
\end{gather}
cf.~equation~(3.73) in~\cite{GLFII}. Therefore, \eqref{Rreal} is clear from \eqref{cRM} and interpolation. As a corollary, this yields reality of $\cR$ and $F$ for real parameters and variables, cf.~\eqref{cRr} and~~\eqref{FR}. (Alternatively, this reality property of the $\cR$-function follows from that of the $R$-function proved in Lemma~2.1 of~III.)

